% part: intuitionistic-logic
% chapter: propositions-as-types
% section: proof-terms

\documentclass[../../../include/open-logic-section]{subfiles}

\begin{document}

\olfileid[id]{int}{pty}{ter}

\olsection{Term Bukti}

Kita memberi asumsi-asumsi dalam \Log{N1i} label variabel (misalnya,
$!A^x$) sebagai label pengosongan. Dalam \Log{N2i}, setiap sekuen
$\Gamma \Sequent !A$ mempunyai suatu daftar pasangan
variabel--!!{formula} dalam konteks~$\Gamma$ di sebelah kiri. Jadi, dalam
\Log{N2i}, setiap sekuen dalam !!a{proof} tidak hanya membawa informasi
tentang apa yang dicapai !!{proof} tersebut sejauh ini (yakni, bahwa
$!A$ telah di-!!{prove}), tetapi juga asumsi mana yang masih terbuka pada setiap
titik (yakni, $\Gamma$). Bagaimana jika kita juga menyertakan dalam
setiap sekuen informasi tentang \emph{bagaimana} $!A$ di-!!{prove}? Untuk
itu, kita dapat mendefinisikan sistem \Log{tN2i} \emph{berlabel term},
yang sekuennya berbentuk $\Gamma \Sequent N : !A$. $\Gamma$ dan $!A$
sama seperti sebelumnya: yang pertama adalah himpunan pasangan berbentuk
$x:!B$, sedangkan yang kedua adalah !!{formula} yang telah ditunjukkan oleh
sub-!!{proof} yang berakhir pada sekuen tersebut sebagai konsekuensi
asumsi-asumsi dalam~$\Gamma$. $N$ adalah term tertentu yang mengodekan
!!{proof} yang berakhir pada sekuen tersebut.

Term yang kita berikan kepada setiap sekuen dibangun dari variabel $x$,
$y$, \dots, dengan menggunakan simbol fungsi yang mengodekan inferensi
yang telah dilaksanakan. Kita memerlukan simbol fungsi tersendiri bagi
setiap aturan inferensi. Simbol fungsi yang bersesuaian dengan aturan
!!{introduction} disebut ``konstruktor,'' sedangkan yang bersesuaian
dengan aturan !!{elimination} disebut ``destruktor.'' Setiap simbol
fungsi menerima sebanyak argumen sebanyak jumlah premis aturannya.
Misalnya, simbol fungsi bagi \Intro{\land} dan \Elim{\land}[1] dapat
berupa $c^\land$ (berargumen dua) dan $d^\land_1$ (berargumen satu). Dalam
\Log{tN2i}, aturan-aturan ini kemudian menjadi
\[
\Axiom$\Gamma_1 \fCenter N: !A$
\Axiom$\Gamma_2 \fCenter M: !B$
\RightLabel{\Intro{\land}}
\BinaryInf$\Gamma_1, \Gamma_2 \fCenter c^\land(N, M) : !A \land !B $
\DisplayProof
\quad
\Axiom$\Gamma \fCenter N: !A \land !B$
\RightLabel{\Elim{\land}[1]}
\UnaryInf$\Gamma \fCenter d^\land_1(N) : !A$
\DisplayProof
\]
Jika suatu aturan inferensi mengosongkan asumsi, yakni aturan itu mempunyai
premis yang menunjukkan !!{formula} berlabel~$x:!A$ dalam suatu premis
tetapi label tersebut tidak ada dalam konteks kesimpulan, kita juga harus
menunjukkan hal ini dalam term bukti. Simbol fungsi yang terkait dengan
aturan semacam itu \emph{mengikat} variabel yang bersesuaian. Misalnya,
aturan \Intro{\lif} beralih dari premis $x:!A, \Gamma \Sequent !B$ menuju
$\Gamma \Sequent !A \lif !B$. Konstruktor $c^{\lif}$ menerima satu
argumen, tetapi juga harus menunjukkan bahwa variabel $x$ menjadi
terikat. Inilah cara term bukti menyatakan bahwa asumsi berlabel~$x$
telah dikosongkan. Misalnya, jika term bukti untuk premis ialah~$N$, kita
dapat menulis term bukti untuk kesimpulan sebagai $c^{\lif}([x]N)$ dan
aturannya sebagai
\[
\Axiom$x:!A, \Gamma \fCenter N: !B$
\RightLabel{\Intro{\lif}}
\UnaryInf$\Gamma \fCenter c^{\lif}([x]N) : !A \lif !B$
\DisplayProof
\]

Agar !!a{proof} dapat direkonstruksi sepenuhnya, kita memerlukan beberapa
informasi tambahan. Jika !!{formula}s premis suatu aturan tidak sepenuhnya
menentukan !!{formula} kesimpulannya, kita harus menambahkan informasi itu
kepada term. Hal ini terjadi dalam \Intro{\lif}, sehingga kita akan
menulis $c^{\lif}([x:!A]N)$. Hal ini juga terjadi dalam \Intro{\lor}
dan~\FalseInt, sehingga kita memakai simbol fungsi yang membawa informasi
tersebut, misalnya,
\[
\Axiom$\Gamma \fCenter N:!A$
\RightLabel{\Intro{\lor}[1]}
\UnaryInf$\Gamma \fCenter c^{\lor{!B}}_1(N):!A \lor !B$
\DisplayProof
\quad
\Axiom$\Gamma \fCenter N:\lfalse$
\RightLabel{\FalseInt}
\UnaryInf$\Gamma \fCenter d^\lfalse_{!A}(N):!A$
\DisplayProof
\]

Dengan gagasan ini, kita dapat membangun sistem deduksi alami bergaya
sekuen yang setiap sekuennya berbentuk $\Gamma \Sequent N : !A$, dengan
$N$ suatu \emph{term bukti}.

Ternyata terdapat hubungan erat antara term bukti semacam itu dan term
dalam apa yang disebut \emph{kalkulus lambda bertipe}. Untuk menghormati
hubungan ini, kita akan memakai simbol yang digunakan dalam literatur
tentang kalkulus lambda, bukan simbol konstruktor dan destruktor umum di
atas. Misalnya, alih-alih menulis $c^{\lif}([x:!A]N)$ kita menulis
$\lambd[\typeof{x}{!A}][N]$; alih-alih $d^{\lif}(N,M)$ kita menulis
$(NM)$; kita menulis $\pair{N}{M}$ alih-alih $c^\land(N,M)$, dan
$\proj{i}{N}$ alih-alih~$d^\land_i(N)$.

\begin{defn}
  Term bukti dihasilkan secara induktif oleh aturan-aturan berikut:
  \begin{enumerate}
  \item Setiap variabel $x$ merupakan term bukti (Aksioma).
  \item Jika $x$ adalah variabel, $!A$ adalah !!a{formula}, dan $N$
    merupakan term bukti, maka $\lambd[\typeof{x}{!A}][N]$ juga merupakan
    term bukti~(\Intro\lif).
  \item Jika $N$ dan $M$ merupakan term bukti, maka $(NM)$ juga merupakan
    term bukti~(\Elim\lif).
  \item Jika $N$ dan $M$ merupakan term bukti, maka $\pair{N}{M}$ merupakan
    term bukti~(\Intro\land).
  \item Jika $N$ merupakan term bukti, maka $\proj{i}{N}$ juga merupakan
    term bukti, dengan $i$ sama dengan $1$ atau~$2$~(\Elim\land[i]).
  \item Jika $N$ merupakan term bukti dan $!A$ merupakan formula, maka
    $\inj[!A]{i}{N}$ merupakan term bukti, dengan $i$ sama dengan $1$
    atau~$2$~(\Intro\lor[i]).
  \item Jika $M$, $N_1$, dan $N_2$ merupakan term bukti, serta $x_1$ dan
    $x_2$ merupakan variabel, maka
    $\dcase{M}{x_1}{N_1}{x_2}{N_2}$ merupakan term
    bukti~(\Elim\lor).
  \item Jika $N$ merupakan term bukti dan $!A$ adalah !!a{formula}, maka
    $\abort{!A}{N}$ merupakan term bukti~(\FalseInt).
  \end{enumerate}
\end{defn}

Tidak setiap term bukti merupakan term bukti dari !!a{proof}. Misalnya,
term $\pair{N}{M}$ mengodekan bukti yang berakhir dengan inferensi
\Intro{\land}, sehingga !!{formula} akhirnya adalah konjungsi
$!A \land !B$. Formula itu tidak dapat menjadi premis mayor suatu aturan
\Elim{\lif}, sehingga $(\pair{N}{M}P)$ tidak dapat menjadi term bukti
dari bukti yang benar. Hanya term bukti yang mengikuti aturan
\Log{tN2ip} yang merupakan term bukti \emph{benar}. Aturan-aturan ini
diberikan dalam \olref{tab:tN2ip}.

\subfile{rules-tN2}

\end{document}
